\subsection{Penjumlahan Vektor dalam \texorpdfstring{$\mathbb{R}^n$}{Rn}}

Penjumlahan vektor dalam $\mathbb{R}^n$ didefinisikan sebagai berikut.  \index{vektor!penjumlahan}

\begin{definition}{Penjumlahan Vektor dalam $\mathbb{R}^n$}\label{def:vectoraddition}
Jika $\vect{u}=\leftB \begin{array}{c}
u_{1} \\
\vdots \\
u_{n}
\end{array}
\rightB,\; \vect{v}= \leftB \begin{array}{c}
v_{1} \\
\vdots \\
v_{n}
\end{array}
\rightB \in \mathbb{R}^{n}$
maka $\vect{u}+\vect{v}\in \mathbb{R}^{n}$ dan didefinisikan oleh
\begin{eqnarray*}
\vect{u}+\vect{v} &=& \leftB \begin{array}{c}
u_{1} \\
\vdots \\
u_{n}
\end{array}
\rightB +  \leftB \begin{array}{c}
v_{1} \\
\vdots \\
v_{n}
\end{array} \rightB\\
& =& \leftB \begin{array}{c}
u_{1}+v_{1} \\
\vdots \\
u_{n}+v_{n}
\end{array}
\rightB
\end{eqnarray*}
\end{definition}

Untuk menjumlahkan vektor, kita cukup menjumlahkan komponen-komponen yang bersesuaian. Oleh karena itu, agar dapat dijumlahkan, vektor-vektor tersebut harus
berukuran sama.

Penjumlahan vektor memenuhi beberapa sifat penting yang diuraikan dalam teorema berikut.

\begin{theorem}{Sifat-Sifat Penjumlahan Vektor}\label{thm:propertiesvectoraddition}
Sifat-sifat berikut berlaku untuk vektor $\vect{u},\vect{v}, \vect{w} \in \mathbb{R}^{n}$.
\begin{itemize}
\item Hukum Komutatif Penjumlahan
\begin{equation*}
\vect{u}+\vect{v}=\vect{v}+\vect{u}
\end{equation*}
\item Hukum Asosiatif Penjumlahan
\begin{equation*}
\left( \vect{u}+\vect{v}\right) +\vect{w}=\vect{u}+\left( \vect{v}+\vect{w}\right)
\end{equation*}
\item Keberadaan Identitas Aditif
\begin{equation}
\vect{u}+\vect{0}=\vect{u}
\label{vectoridentity}
\end{equation}
\item Keberadaan Invers Aditif
\begin{equation*}
\vect{u}+\left( -\vect{u}\right) =\vect{0}
\end{equation*}
\end{itemize}
\end{theorem}

Identitas aditif yang ditunjukkan dalam persamaan
\ref{vectoridentity} juga disebut \textbf{vektor nol}\index{vektor nol},
yaitu vektor $n \times 1$ yang semua komponennya sama dengan $0$.
Lebih lanjut, $-\vect{u}$ hanyalah vektor yang semua komponennya memiliki
nilai yang sama dengan komponen-komponen $\vect{u}$ tetapi bertanda kebalikan; ini tidak lain adalah
$(-1)\vect{u}$. Hal ini akan dibuat lebih eksplisit pada bagian berikutnya
ketika kita mengkaji perkalian skalar pada vektor. Perhatikan bahwa pengurangan didefinisikan sebagai $\vect{u}-\vect{v} = \vect{u}+\left(
-\vect{v} \right)$ \index{vektor!pengurangan}.
